\subsection{NOTATION AND TERMINOLOGY}

The definitions to be given in the remainder of this section apply to an arbitrary order (or preorder) relation R\{x, y\} between x and y, but will be used mainly in the case where R\{x, y\} is written \(x \leqslant y\) *(by analogy with the usual order relation between integers or real numbers)* (or \(x \subset y\) , or by means of some analogous sign); we shall therefore state them only for the notation \(x \leqslant y\) , and leave to the reader the task of transcribing them into other notations. When R\{x, y\} is written \(x \leqslant y\) , then \(y \geqslant x\) is synonymous with \(x \leqslant y\) , and these relations are read ``x is smaller than y'', or ``x is less than y'', or ``y is larger than x'' or ``y is greater than x''. The relation \(x \geqslant y\) is then the preorder relation (between x and y) opposite to \(x \leqslant y\) .

By abuse of language, we shall often speak of ``the relation \(\leqslant\) '' instead of ``the relation \(x \leqslant y\) ''; in this case ``the relation \(\geqslant\) '' is the opposite of ``the relation \(\leqslant\) ''. We remark also that, in the same proof, we may often use the same sign \(\leqslant\) to denote several different order relations when there is no risk of confusion.

The conditions for a relation written \(x \leqslant y\) to be an order relation on a set \(\mathbf{E}\) are as follows:

\((\mathbf{RO}_{\mathbf{I}})\) The relation '' \(x \leqslant y\) and \(y \leqslant z\) '' implies \(x \leqslant z\) .

\((\mathbf{RO}_{\mathbf{II}})\) The relation '' \(x \leqslant y\) and \(y \leqslant x\) '' implies \(x = y\) .

\((\mathbf{RO}_{\mathbf{III}})\) The relation \(x \leqslant y\) implies `` \(x \leqslant x\) and \(y \leqslant y\) ''.

\((\mathbf{RO}_{\mathbf{IV}})\) The relation \(x \leqslant x\) is equivalent to \(x \in \mathbf{E}\) .

If we leave out condition \((\mathrm{RO}_{\mathbf{II}})\) , we have the conditions for \(x \leqslant y\) to be a preorder relation on E.

When an order relation is written \(x \leqslant y\) we shall write x \textless{} y (or y \textgreater{} x) for the relation `` \(x \leqslant y\) and \(x \neq y\) ''; these relations are read ``x is strictly smaller than y'', or ``x is strictly less than y'', or ``y is strictly larger than x'', or ``y is strictly greater than x''.

The example of the inclusion relation shows that the negation of \(x \leqslant y\) (sometimes denoted by \(x \nleqslant y\) ) is not necessarily equivalent to \(y < x\) (cf.~no. 12).

C58. Let \(\leqslant\) be an order relation, and let x, y be two distinct letters. The relation \(x \leqslant y\) is equivalent to ``x \textless{} y or x = y''. Each of the relations `` \(x \leqslant y\) and \(y < z\) '', `` \(x < y\) and \(y \leqslant z\) '' implies x \textless{} z.

The first assertion follows from the criterion \(\mathbf{A} \Longrightarrow ((\mathbf{A}\) and (not \(\mathbf{B}\) )) or \(\mathbf{B}\) (Chapter I, §3, criterion C24). To prove the second assertion, we remark that each of the hypotheses implies \(x \leqslant z\) , by transitivity; and the relation ( \(x = z\) and \(x \leqslant y\) and \(y \leqslant z\) ) would imply \(x = y = z\) , contrary to the hypothesis.

In order to make matters easier and to replace metamathematical criteria by mathematical theorems we shall usually consider a theory \(\mathfrak{G}\) which contains the axioms and axiom schemes of the theory of sets, and in addition, two constants \(\mathbf{E}\) and \(\Gamma\) satisfying the axiom

\textbf{``Γ is an ordering on the set E'' (no. 1).}

We shall denote by \(x \leqslant y\) the relation \(y \in \Gamma \langle x \rangle\) , and we shall say that the set \(E\) is ordered by the ordering \(\Gamma\) (or by the order relation \(y \in \Gamma \langle x \rangle\) ) (cf.~Chapter IV, §1).

¶ Whenever, in \(\mathfrak{G}\) , \(\Gamma\) is a preordering on \(\mathbf{E}\) , we say likewise that \(\mathbf{E}\) is preordered by the preordering \(\Gamma\) .

In some situations (for example in the following definition) the theories which we shall consider are a little more complicated. We shall leave it to the reader to make explicit the constants and axioms of such theories.

Let E, \(E'\) be two sets ordered by orderings \(\Gamma\) , \(\Gamma'\) . An isomorphism of E onto \(E'\) (for the orderings \(\Gamma\) and \(\Gamma'\) ) is a bijection f of E onto \(E'\) such that the relations \(x \leqslant y\) and \(f(x) \leqslant f(y)\) are equivalent (cf.~Chapter IV, §1).

